%\documentclass[11pt]{book}
\documentclass[a4paper,11pt]{article}
\usepackage{graphicx}
\usepackage{pifont}
\usepackage{ulem}
\usepackage{multirow}
\usepackage{supertabular}
\usepackage{fancyhdr}
\usepackage{rotating}
\usepackage{verbatim}
\usepackage{appendix}
\usepackage{placeins}
\usepackage[frenchb]{babel}

\voffset -1.5cm
\hoffset 0.0cm
\textheight 25cm
\textwidth 16cm
\topmargin 0.0cm
\oddsidemargin 0.0cm
\evensidemargin 0.0cm

\newcommand{\VROT}[1]{\begin{turn}{90}#1\end{turn}}     % vertical text
\def\X{$\surd$}     % tick symbol

\usepackage{hyperref}%                   % Utilisation de HyperTeX

\usepackage[utf8x]{inputenc}

%\lhead{\includegraphics[width=1.0cm]{logo_unisim.pdf}}

\hypersetup{
    pdfauthor   = {Julien Lisita},%
    pdftitle    = {Mémoire de stage},%
    pdfsubject  = {Développement d'un simulateur de jeu d'instructions},%
    pdfkeywords = {Environnement de simulation, exploration d'architecture, conception des systèmes embarqués, vérification du système},%
    pdfcreator  = {rubber},%
    pdfproducer = {rubber}
}


%\geometry{hmargin=3.0cm,vmargin=2.5cm}

\begin{document}
\selectlanguage{french}
\setlength{\parskip}{0.5em}%0.15cm}
\baselineskip 14pt


\title{\bf \huge Développement d'un simulateur de jeu d'instructions}
\author{Julien Lisita}
\date{1$^{er}$ Septembre 2014}

\maketitle

\newpage

% \begin{titlepage}
% %\frontmatter
% \begin{figure}
% \begin{center}
% %\includegraphics[width=8cm,natwidth=11.0cm,natheight=11.0cm]{logo_unisim.pdf}
% \maketitle
% \end{center}
% \end{figure}
% \end{titlepage}
% %\input{abstract}
% %\input{acknowledgements}
% %\tableofcontents


%\addcontentsline{toc}{chapter}{Table des mati�res} 
\tableofcontents \markboth{Table of content}{Table of content}
\listoffigures

%\mainmatter

\section{Introduction}
\label{introduction}
\section{Introduction}
\label{sec:str7_validation_introduction}

In the developpement of new simulators an important task is the validation of the simulator against the real target architecture.
The objective of this task is to check that the behavior of the simulator closely matches that of the target architecture. 
While a complete validation of a simulator is probably an impossible task, or at least very time consuming to be affordable, we present in this document the validation we have used over the STR7 microcontroller simulator and its results. 
We consider this validation to be enough to start using the simulator in the industrial and research domains with confidence.

This document shows the different aspects of the simulator that have been validated.
Section~\ref{sec:str7_validation_description} and Section~\ref{sec:str7_validation_tools} briefly discuss the simulated architecture and the tools used during the validation.
Section~\ref{sec:str7_validation_user_level} and Section~\ref{sec:str7_validation_system_level} introduce the different validation tests that have been performed on the main component of the STR7 simulator, that is, the ARM7TDMI processor.
More concretely Section~\ref{sec:str7_validation_user_level} shows the tests that were performed at the user level, and Section~\ref{sec:str7_validation_system_level} shows the tests performed at the system level.
Finally, in Section~\ref{sec:str7_validation_microcontroller_validation} discusses the validation of the complete microcontroller in an industrial environment.


\section{Contexte du stage}
\section{Prise en main du problème}
\section{Méthodologie}
\section{Évaluation}
\section{Discussion des résultats}
\section{Conclusion}
\section{Apport du stage}

\begin{appendix}
\section{Implantation et couverture du jeu d'instructions par benchmarking}
\begin{center}
\tablehead{\hline
\multicolumn{1}{|c|}{\textbf{Instruction}} &
\VROT{Inplantée} &
\VROT{Basicmath (large)} &
\VROT{Basicmath (small)} &
\VROT{Bitcount (large)} &
\VROT{Bitcount (small)} &
\VROT{Dijkstra (large)} &
\VROT{Dijkstra (small)} &
\VROT{Stringsearch (large)} &
\VROT{Stringsearch (small)} &
\VROT{Blowfish (large)} &
\VROT{Blowfish (small)} &
\VROT{Rijndael (large)} &
\VROT{Rijndael (small)} &
\VROT{SHA (large)} &
\VROT{SHA (large)} &
\VROT{adpcm (large)} &
\VROT{adpcm (small)} &
\VROT{CRC32 (large)} &
\VROT{CRC32 (small)}\\
\hline}
\tabletail{\hline}
\begin{supertabular}{|p{4.0cm}|p{0.25cm}|p{0.25cm}|p{0.25cm}|p{0.25cm}|p{0.25cm}|p{0.25cm}|p{0.25cm}|p{0.25cm}|p{0.25cm}|p{0.25cm}|p{0.25cm}|p{0.25cm}|p{0.25cm}|p{0.25cm}|p{0.25cm}|p{0.25cm}|p{0.25cm}|p{0.25cm}|p{0.25cm}|}
\multicolumn{20}{|c|}{\textbf{add}}\\
\hline
\multicolumn{1}{|p{4.0cm}|}{\scriptsize \texttt{add Rd, Rs}} &
\multicolumn{1}{p{0.25cm}|}{\X} &
\multicolumn{1}{p{0.25cm}|}{} &
\multicolumn{1}{p{0.25cm}|}{} &
\multicolumn{1}{p{0.25cm}|}{} &
\multicolumn{1}{p{0.25cm}|}{} &
\multicolumn{1}{p{0.25cm}|}{} &
\multicolumn{1}{p{0.25cm}|}{} &
\multicolumn{1}{p{0.25cm}|}{} &
\multicolumn{1}{p{0.25cm}|}{} &
\multicolumn{1}{p{0.25cm}|}{} &
\multicolumn{1}{p{0.25cm}|}{} &
\multicolumn{1}{p{0.25cm}|}{} &
\multicolumn{1}{p{0.25cm}|}{} &
\multicolumn{1}{p{0.25cm}|}{} &
\multicolumn{1}{p{0.25cm}|}{} &
\multicolumn{1}{p{0.25cm}|}{} &
\multicolumn{1}{p{0.25cm}|}{} &
\multicolumn{1}{p{0.25cm}|}{} &
\multicolumn{1}{p{0.25cm}|}{}\\
\hline
\multicolumn{1}{|p{4.0cm}|}{\scriptsize \texttt{add Rd, Rx, Ry << sa}} &
\multicolumn{1}{p{0.25cm}|}{\X} &
\multicolumn{1}{p{0.25cm}|}{} &
\multicolumn{1}{p{0.25cm}|}{} &
\multicolumn{1}{p{0.25cm}|}{} &
\multicolumn{1}{p{0.25cm}|}{} &
\multicolumn{1}{p{0.25cm}|}{} &
\multicolumn{1}{p{0.25cm}|}{} &
\multicolumn{1}{p{0.25cm}|}{} &
\multicolumn{1}{p{0.25cm}|}{} &
\multicolumn{1}{p{0.25cm}|}{} &
\multicolumn{1}{p{0.25cm}|}{} &
\multicolumn{1}{p{0.25cm}|}{} &
\multicolumn{1}{p{0.25cm}|}{} &
\multicolumn{1}{p{0.25cm}|}{} &
\multicolumn{1}{p{0.25cm}|}{} &
\multicolumn{1}{p{0.25cm}|}{} &
\multicolumn{1}{p{0.25cm}|}{} &
\multicolumn{1}{p{0.25cm}|}{} &
\multicolumn{1}{p{0.25cm}|}{}\\
\hline
\multicolumn{20}{|c|}{\textbf{add\{cond4\}}}\\
\hline
\multicolumn{1}{|p{4.0cm}|}{\scriptsize \texttt{add\{cond4\} Rd, Rx, Ry}} &
\multicolumn{1}{p{0.25cm}|}{\X} &
\multicolumn{1}{p{0.25cm}|}{} &
\multicolumn{1}{p{0.25cm}|}{} &
\multicolumn{1}{p{0.25cm}|}{} &
\multicolumn{1}{p{0.25cm}|}{} &
\multicolumn{1}{p{0.25cm}|}{} &
\multicolumn{1}{p{0.25cm}|}{} &
\multicolumn{1}{p{0.25cm}|}{} &
\multicolumn{1}{p{0.25cm}|}{} &
\multicolumn{1}{p{0.25cm}|}{} &
\multicolumn{1}{p{0.25cm}|}{} &
\multicolumn{1}{p{0.25cm}|}{} &
\multicolumn{1}{p{0.25cm}|}{} &
\multicolumn{1}{p{0.25cm}|}{} &
\multicolumn{1}{p{0.25cm}|}{} &
\multicolumn{1}{p{0.25cm}|}{} &
\multicolumn{1}{p{0.25cm}|}{} &
\multicolumn{1}{p{0.25cm}|}{} &
\multicolumn{1}{p{0.25cm}|}{}\\
\hline
\end{supertabular}
\end{center}
 
\end{appendix}

\bibliographystyle{plain}
\bibliography{memoire_lisita/biblio}

\end{document}
